\scanpage{34}%
\mspage{19}%
laire s'écrit
\[
 b_2^{-n}\,P(b_2,\Delta) \qquad \text{pour }P\text{ isobare}\quad ,
\]
et que cette forme modulaire est définie aussi pour $j=0$ si et seulement si elle s'écrit comme polyn\mbox{\begingroup\edef\encoutersep{\the\dimexpr.5\pgflinewidth\relax}\setbox0=\hbox{\mbox{$\overset{\inserted{\handwritten{\wedge}}}{\mathrm{o}}$}}\edef\encwda{\the\wd0}\edef\enchta{\the\ht0}\edef\encdpa{\the\dp0}\rlap{\smash{\tikz[baseline=0pt]{\begin{scope}[shift={(\the\dimexpr \encwda/2\relax,\the\dimexpr -\encdpa/2+\enchta/2\relax)}]\coordinate (enc-oend-center) at (0pt,0pt);\handwritten{\draw[thin,gray,rotate around={-20:(enc-oend-center)}] (enc-oend-center) ellipse (0.64em and 1.15em);}\end{scope}\pgfresetboundingbox\path[use as bounding box] (0pt,\the\dimexpr-\encdpa\relax) rectangle (\encwda,\enchta);}}}\copy0\endgroup}me en $b_2$ et $\Delta$ .

\handwritten{Interpr\unsure{é|e}tation de $a_1$, $b_2$ comme forme de Hasse-Witt ?}

\handwritten{Formes modulaires \unsure{sur|en} un corps quelconque ? Sur un ann\unsure{.|,} de val. discrète\ldots}

\handwritten{Hasse-Witt sur un corps de car $p>0$ quelconque ?}

\noindent\handwritten{\tikz{\draw[thin] (0,0) -- (3em,.06em);}}

\[
\handwritten{\left\{\begin{aligned}
 &\text{Relèvement canonique de Serre par formules universelles}\,\big|\\
 &\text{Loi de groupe formel sur }{\mathcal{P}}^{\infty}_{\unsure{M|\mu}/S}\ \text{: formules explicites\ldots}
\end{aligned}\right.}
\]
% archive sheet number: 33
